\section*{Chapitre \ref{ch:b3:orogenic-wedges} --- Chaînes de chevauchement et prismes orogéniques}

\begin{solution}{exo:b3:orogenic-wedges:1}
$\sin\ang{30} = 0.5$~: $K = 2\times0.5/0.5 = 2$ et $\sigma'_1/\sigma'_3 = 1 +
K = 3$.
\end{solution}

\begin{solution}{exo:b3:orogenic-wedges:2}
$\phi = \arctan0.6 = \ang{31.0}$, $\sin\phi = 0.514$, $K = 2.12$. Sur une
base horizontale, $\alpha = 0.3/3.12 = \qty{0.0962}{rad} = \ang{5.5}$. Sur une
base inclinée de $\ang{3} = \qty{0.0524}{rad}$, $\alpha + \beta = 0.352/3.12 =
\qty{0.113}{rad} = \ang{6.5}$, donc $\alpha = \ang{3.5}$.
\end{solution}

\begin{solution}{exo:b3:orogenic-wedges:3}
Profondeur $= (T - 10)/30$~: apatite (U--Th)/He de 1.3 à \qty{2.3}{km}~; apatite à
traces de fission de 2.7 à \qty{3.7}{km}~; zircon à traces de fission $270/30 = \qty{9.0}{km}$. À
\qty{5}{km}, la roche se trouvait à $10 + 150 = \qty{160}{\celsius}$~: zircon
(U--Th)/He, qui se ferme entre 160 et \qty{200}{\celsius}, à 5.0 à
\qty{6.3}{km} de profondeur.
\end{solution}

\begin{solution}{exo:b3:orogenic-wedges:4}
Longueur horizontale $2/\tan\ang{30} = \qty{3.46}{km}$~; le long du pendage
$2/\sin\ang{30} = \qty{4.0}{km}$. La couche est soulevée de la hauteur de la \omterm{def:b3:orogenic-wedges:thrust-geometry}{rampe},
\qty{2}{km}. Le point de coupure doit parcourir toute la \omterm{def:b3:orogenic-wedges:thrust-geometry}{rampe}~: un glissement de
\qty{4.0}{km}.
\end{solution}

\begin{solution}{exo:b3:orogenic-wedges:5}
$\phi = \arctan1.03 = \ang{45.8}$, $K = 5.08$~; $\beta = \qty{0.105}{rad}$.
Pour $\lambda = 0.4$~: $\alpha + \beta = (0.105 + 0.85\times0.6)/(1 +
5.08\times0.6) = 0.615/4.05 = \qty{0.152}{rad} = \ang{8.7}$, $\alpha =
\ang{2.7}$. Pour $\lambda = 0.9$~: $0.190/1.51 = \qty{0.126}{rad} = \ang{7.2}$,
$\alpha = \ang{1.2}$. Augmenter la pression de fluide interstitiel abaisse le \omterm{def:b3:orogenic-wedges:critical-taper}{biseau critique}~:
le prisme devient surcritique, glisse vers l'avant sur sa base sans
épaissir, et son front progresse jusqu'à ce que l'accrétion et l'érosion le ramènent
au nouveau \omterm{def:b3:orogenic-wedges:critical-taper}{biseau critique}.
\end{solution}

\begin{solution}{exo:b3:orogenic-wedges:6}
$S = A/r = 6.0/1.5 = \qty{4.0}{km}$. Interprétée comme un \omterm{def:b3:orogenic-wedges:thrust-folds}{pli de décollement}, $S = A/h =
6.0/3.0 = \qty{2.0}{km}$~: la moitié de la valeur réelle. Les deux interprétations doivent être
distinguées par la forme du pli et la présence d'une \omterm{def:b3:orogenic-wedges:thrust-geometry}{rampe}.
\end{solution}

\begin{solution}{exo:b3:orogenic-wedges:7}
$K = 2.12$ et $1 + K(1 - 0.4) = 2.27$. L'équation donne $\mu_b(1 -
\lambda_b) = (\alpha + \beta)\times2.27 - \beta$, avec $\beta =
\qty{0.0349}{rad}$. Jura~: $0.0611\times2.27 - 0.0349 = 0.10$. Sud~:
$0.218\times2.27 - 0.0349 = 0.46$. Byerlee à pression hydrostatique~:
$0.85\times0.6 = 0.51$. La base au sud est proche de celle de Byerlee~; la
base du Jura est presque cinq fois plus faible, comme attendu pour un décollement dans
halite, dont le frottement est faible, sous des fluides non surpressurisés.
\end{solution}

\begin{solution}{exo:b3:orogenic-wedges:8}
$\kappa = \qty{e-6}{m^2/s} = \qty{31.6}{km^2/Ma}$, et $v =
\qty{1}{km/Ma}$~: $\mathrm{Pe} = 30/31.6 = 0.951$. Le facteur $0.951/(1 -
e^{-0.951}) = 1.55$ fait passer le gradient de 30 à \qty{46.5}{\celsius/km}.
Pour \qty{65}{\celsius}~: $1 - e^{-\mathrm{Pe}\,z/L} = (55/900)(1 - e^{-0.951})
= 0.0375$, donc $\mathrm{Pe}\,z/L = 0.0382$ et $z = \qty{1.21}{km}$, contre
$55/30 = \qty{1.83}{km}$ sans érosion.
\end{solution}

\begin{solution}{exo:b3:orogenic-wedges:9}
$\mathrm dh/\mathrm dt = \qty{0.8}{km}/\qty{1.2}{Ma} = \qty{0.67}{mm/an}$.
Avec $f = 0.5$, $v = 0.67\times(1 - 0.5) = \qty{0.33}{mm/an}$.
\end{solution}

\begin{solution}{exo:b3:orogenic-wedges:10}
Avec $\alpha = 0$ et $u = 1 - \lambda$~: $\beta(1 + Ku) = \beta + \mu_b u$,
donc $\beta Ku = \mu_b u$ et $\beta^* = \mu_b/K$, quel que soit $u$. Pour Taïwan,
$0.85/5.08 = \qty{0.167}{rad} = \ang{9.6}$. Sur une base plus pentue, la
pente critique de surface est négative~: tout prisme à surface horizontale ou inclinée vers l'avant-pays
est déjà surcritique et glisse sur sa base sans s'épaissir. Une
chaîne de chevauchement ne peut pas y acquérir sa forme normale.
\end{solution}

\begin{solution}{exo:b3:orogenic-wedges:11}
Comme $\Lambda \to \infty$ et $v \to 0$, $m \to 0$, donc $e^{-mz_c} \to 1$, et
le géotherme est rectiligne, $G_0 = G(z_c) = G$~: $f \to (G - \Gamma)/G = 1 -
\Gamma/G = 1 - 6.5/30 = 0.78$. La pente vaut alors $v/(1 - f) = vG/\Gamma =
4.6v$. Pour $v = \qty{1}{mm/an}$, $v/2\kappa = \qty{0.0158}{km^{-1}}$. À
\qty{5}{km}~: $m = 0.0158 + \sqrt{0.0158^2 + 1.257^2} = \qty{1.273}{km^{-1}}$,
$e^{-1.273\times1.21} = 0.215$, $f = 40.0\times0.215/44.7 = 0.19$. À
\qty{20}{km}~: $m = 0.0158 + 0.3146 = \qty{0.330}{km^{-1}}$, $e^{-0.399} =
0.671$, $f = 40.0\times0.671/44.7 = 0.60$.
\end{solution}

\begin{solution}{exo:b3:orogenic-wedges:12}
La coupe est un triangle de base $W$ et de hauteur $W\tan(\alpha + \beta)$~:
$A = W^2\tan(\alpha + \beta)/2$. Avec $\mathrm dA/\mathrm dt = F_a$ et $A =
0$ à $t = 0$, $A = F_at$, ce qui donne $W$. Avec $\tan\ang{6} = 0.105$~:
$W = \sqrt{2\times40\times2/0.105} = \qty{39}{km}$ après \qty{2}{Ma}~; $t =
60^2\times0.105/(2\times40) = \qty{4.7}{Ma}$ pour atteindre \qty{60}{km}.
\end{solution}

\begin{solution}{pb:b3:orogenic-wedges:1}
\textbf{1.} $4.5 - 1.5 = \qty{3.0}{km}$ de hauteur~; $3.0/\tan\ang{25} =
\qty{6.43}{km}$ de longueur horizontale~; $3.0/\sin\ang{25} = \qty{7.10}{km}$ le long
du pendage.

\textbf{2.} Par la hauteur de la \omterm{def:b3:orogenic-wedges:thrust-geometry}{rampe}, \qty{3.0}{km}.

\textbf{3.} $S = A/r = 16.5/3.0 = \qty{5.5}{km}$.

\textbf{4.} Même valeur, \qty{16.5}{km^2}~: dans un \omterm{def:b3:orogenic-wedges:thrust-folds}{pli de rampe}, la surface en excès
de chaque couche au-dessus du \omterm{def:b3:orogenic-wedges:thrust-geometry}{palier} supérieur vaut $Sr$. Interprétée comme un \omterm{def:b3:orogenic-wedges:thrust-folds}{pli de décollement}, $16.5/4.0 =
\qty{4.1}{km}$~: un quart trop faible.

\textbf{5.} $4\times(6.0 - 2.5) = \qty{14.0}{km}$, soit deux fois et demie le
\qty{5.5}{km} de rejet sur la \omterm{def:b3:orogenic-wedges:thrust-geometry}{rampe} frontale.

\textbf{6.} L'écaille frontale, celle qui est la plus proche de l'avant-pays~: chaque nouvelle \omterm{def:b3:orogenic-wedges:thrust-geometry}{rampe}
se forme dans le mur, en avant de la précédente. Chaque écaille, portée sur sa \omterm{def:b3:orogenic-wedges:thrust-geometry}{rampe},
soulève le chevauchement de toit et la couverture qui le surmonte, comme un \omterm{def:b3:orogenic-wedges:thrust-folds}{pli de rampe} soulève son
toit, si bien que ce chevauchement est plissé au-dessus de l'empilement.

\textbf{7.} $\alpha = \arctan(2.1/60) = \ang{2.00}$; $\beta =
\arctan(4.2/60) = \ang{4.00}$~; biseau \ang{6.0}.

\textbf{8.} $2.4 + 8.2 = \qty{10.6}{km}$, inférieur à environ \qty{15}{km}~: la
base est encore frictionnelle, et la théorie de Coulomb s'applique.

\textbf{9.} $\phi = \ang{45.8}$, $\sin\phi = 0.717$, $K = 5.08$. À sec~:
$\alpha + \beta = (0.0699 + 0.85)/6.08 = \qty{0.151}{rad} = \ang{8.7}$,
$\alpha = \ang{4.7}$.

\textbf{10.} $\lambda = 1000/2500 = 0.40$: $(0.0699 + 0.51)/4.05 =
\qty{0.143}{rad} = \ang{8.2}$, $\alpha = \ang{4.2}$. La pente observée \ang{2.0}
rendrait le prisme fortement sous-critique, s'épaississant plutôt que d'avancer~:
les fluides doivent être en surpression.

\textbf{11.} Avec $T = \qty{0.1049}{rad}$ et $u = 1 - \lambda$, $T(1 + Ku) =
\beta + \mu_bu$ donne $u = (\beta - T)/(TK - \mu_b) = (0.0699 -
0.1049)/(0.5327 - 0.85) = 0.110$: $\lambda = 0.89$.

\textbf{12.} $\sigma_v = 2500\times9.81\times10\,600 = \qty{260}{MPa}$; $P =
0.89\times260 = \qty{231}{MPa}$~; hydrostatique $1000\times9.81\times10\,600 =
\qty{104}{MPa}$. Les fluides portent \qty{127}{MPa} de plus qu'une colonne de
eau le ferait.

\textbf{13.} $(0.0699 + 0.85\times0.3)/(1 + 5.08\times0.3) = 0.325/2.52 =
\qty{0.129}{rad} = \ang{7.4}$, $\alpha = \ang{3.4}$. Le prisme, à \ang{2.0},
est désormais sous-critique~: il cesse d'avancer et s'épaissit par \omterm{def:b3:orogenic-wedges:thrust-geometry}{chevauchement hors
séquence} à l'arrière de son front, jusqu'à ce que sa pente de surface atteigne \ang{3.4}.

\textbf{14.} $F_a = \qty{10}{km/Ma}\times\qty{4.0}{km} = \qty{40}{km^2/Ma}$,
soit \qty{4.0e4}{m^3/an} par kilomètre de chaîne.

\textbf{15.} $\bar e = 40/60 = \qty{0.67}{km/Ma} = \qty{0.67}{mm/an}$.

\textbf{16.} $40\times100 = \qty{4000}{km^3/Ma} = \qty{4.0e6}{m^3/an}$, ou
$\qty{1.0e10}{kg/an}$ (10 millions de tonnes par an).

\textbf{17.} $\mathrm dW/\mathrm dt = F_a/(W\tan(\alpha + \beta)) =
40/(60\times0.105) = \qty{6.3}{km/Ma}$.

\textbf{18.} Non~: à l'état stationnaire, l'érosion doit encore évacuer tout le
flux entrant, si bien que la vitesse moyenne reste \qty{0.67}{mm/an}. Mais la roche ajoutée
sous l'arrière est soulevée à cet endroit~: l'érosion et l'exhumation seraient les plus rapides
sur l'arrière, là où se trouveraient les âges de refroidissement les plus jeunes.

\textbf{19.} $v = \qty{0.667}{km/Ma}$: $\mathrm{Pe} = 0.667\times30/31.6 =
0.634$; $G_0 = 30\times0.634/(1 - e^{-0.634}) = 30\times1.35 =
\qty{40.5}{\celsius/km}$.

\textbf{20.} $1 - e^{-\mathrm{Pe}\,z/L} = (55/900)(1 - e^{-0.634}) = 0.0287$,
donc $\mathrm{Pe}\,z/L = 0.0291$ et $z_c = \qty{1.38}{km}$~; $G(z_c) =
40.5\,e^{-0.0291} = \qty{39.3}{\celsius/km}$.

\textbf{21.} $1.38/0.667 = \qty{2.07}{Ma}$, l'âge de l'échantillon à
\qty{2.00}{km}~: à l'altitude moyenne, la déflexion de l'isotherme est nulle.

\textbf{22.} $v/2\kappa = 0.667/63.1 = \qty{0.0106}{km^{-1}}$, $2\pi/\Lambda =
\qty{0.314}{km^{-1}}$: $m = 0.0106 + 0.3143 = \qty{0.325}{km^{-1}}$;
$e^{-0.325\times1.38} = 0.639$; $f = (40.5 - 6.5)\times0.639/39.3 = 0.552$.

\textbf{23.} L'altitude moyenne est \qty{2.00}{km} et l'âge moyen
\qty{2.068}{Ma}~; $\sum\Delta t\,\Delta h = 0.823$ et $\sum\Delta t^2 =
0.552$, donc la pente vaut \qty{1.49}{km/Ma}~: une vitesse apparente de
\qty{1.49}{mm/an}.

\textbf{24.} $v = 1.49\times(1 - 0.552) = \qty{0.667}{mm/an}$.

\textbf{25.} Le \omterm{def:b3:orogenic-wedges:thermochronology}{thermochronomètre} et le bilan des flux concordent~: le prisme est
à l'état stationnaire, et prendre la pente apparente au pied de la lettre aurait
surestimé la vitesse de plus du double. La vitesse d'exhumation du prisme
de Kestrel est \textbf{\qty[text-series-to-math=true]{0.67}{mm/an}}.
\end{solution}
